\begin{proposition}\label{prop:four-signed}
One has $\zr(11,4)=\zs(11,4)=29$ and $\zr(12,4)=\zs(12,4)=32$.
\end{proposition}
\begin{proof}
For $m\ge6$, pair counting gives $z(m,4)=m+6$. Equality forces six
rows containing the six distinct column pairs and $m-6$ singleton rows.
The nonincreasing quadruples of anchor multiplicities classify all maximum
bases up to row and column permutations. There are six types at eleven
rows and nine at twelve rows. The respective cell caps are 30 and 33.
Complete signed enumeration excludes both targets:
\begin{center}
\begin{tabular}{lrrr}
\toprule
Grid&Maximum-base types&Target pairs&Search states\\
\midrule
$11\times4$&6&13&189,363,352\\
$12\times4$&9&15&379,946,285\\
\bottomrule
\end{tabular}
\end{center}
Every base is exhausted, with one hole in the eleven-row target and none
in the twelve-row target. The per-base inputs and counts are recorded in S4.

The search branches on a remaining
cell: either it is a hole if the remaining budget permits, or it is paired
with each possible partner. Failed signed subfamilies are excluded by complete-pair heredity.
Deleting a complete displayed form grounds all rectangle deviations incident
with its removed cells, so the signed derivation on retained cells replays.
A matching-component bound rejects only insufficient remaining capacity.
The only symmetry reduction swaps two wholly untouched rows with identical
fixed supports, while fixing the pivot, selected pairs and chosen holes.
Thus every branch omission preserves a representative of every completion.

The independent reader matches the node and depth counts of the archived
signed union-find search. It instead
rebuilds a parity graph, colors its components by breadth-first search,
and grounds components containing a zero, an odd cycle, or opposite selected
inequalities. Only a grounded selected-pair deviation enables its
identification equations. The algorithms use the signed criterion throughout;
a recursive-line rejection is never a pruning condition. The source audit,
complete inputs, calibration cases and termination logs are in packet S4.

The recursive lower values 29 and 32 were already obtained in \cite{CC}.
Explicit twelve- and fourteen-pair witnesses are also replayed by a separate
actual-cell rule reader in S4. The hierarchy gives the asserted equalities.
These are signed upper bounds. The actual-SOS intervals remain
$29\le\zt(11,4)\le30$ and $32\le\zt(12,4)\le33$.
\end{proof}

\begin{corollary}\label{cor:four-exception}
For every $m\ge1$ other than possibly $m=14$,
$\zr(m,4)=\zs(m,4)$.
\end{corollary}
\begin{proof}
At one row every cell is fixed. The exact values for two through eight rows
are in \cite{LQ,XY};
Theorem~\ref{thm:further-finite} covers nine and ten rows, and
Proposition~\ref{prop:four-signed} covers eleven and twelve.
At thirteen rows and every $m\ge15$, the recursive constructions
in \cite[Table~1 and Theorem~1.1]{CC} attain the common cell cap. At fourteen rows, the
remaining interval is $37\le\zr\le\zs\le38$.
\end{proof}


